\section{Prospective empirical programme and falsification rules}\label{si:designs}

The designs below are proposals, not completed tests. No registration identifier, locked protocol, or field intervention is claimed. Each should be registered before outcome inspection, with data provenance, estimand, decision rule, and deviations recorded.

\subsection{Design 1: governance-event panel}

\textbf{Unit and data.} Resource--jurisdiction--assessment cycle. Link dated raw observation, assessment data freeze, advice release, committee decision, legal instrument, implementation, compliance, and ecological response. Record date uncertainty and source document for each event. Separate \(T_{\rm obs}\), \(T_{\rm assess}\), \(\Tr\), \(\tau_{\rm dec}\), \(\tau_{\rm dep}\), and \(\tau_{\rm eco}\); do not impute one ``governance lag'' from another.

\textbf{Estimand.} Distribution of effective responsive-review intervals and the impulse/phase response of command, extraction, and stock to a measured assessment innovation, conditional on operating-model class.

\textbf{Identification.} The panel is descriptive until independent variation or a design-based assumption identifies the effect. Use date ranges where records are incomplete; do not encode a scheduled annual TAC as a 0.5-year response merely because in-season closures exist without trigger evidence.

\textbf{Falsification.} A hypothesised sample-and-hold mechanism is weakened if decisions precede the information they are said to respond to, if implementation does not follow the decision, or if the predicted command--stock phase order fails on held-out cycles.

\subsection{Design 2: quasi-experimental schedule variation}

Identify an external schedule change plausibly independent of current ecological innovations. Candidate approaches include event studies, difference-in-differences, synthetic controls, interrupted time series, and state-space interventions, selected according to the institutional change and comparison structure. Define treatment as a change in responsive decision opportunities, not simply a change in implementation lag. A nominal calendar reform is not a treatment if the same rule is still applied to the same information on the same effective schedule.

\textbf{Estimand.} Change in extraction and ecological outcomes attributable to a change in the review operator, with pre-specified direct and lagged effects.

\textbf{Identification checks.} Pre-trends, concurrent rule and data changes, ecological shocks, serial correlation, spillovers, anticipatory behaviour, implementation/compliance, and measurement-vintage changes. Use placebo dates and placebo stocks where feasible.

\textbf{Falsification.} If estimated schedule variation is explained by an ecological shock or if outcomes move before the schedule change, the causal interpretation is rejected or narrowed.

\subsection{Design 3: out-of-sample mechanism comparison}

Fit competing models with distinct causal orderings: (i) ecological forcing only; (ii) cohort resonance or recruitment structure; (iii) institutional feedback with a pre-specified rule sign and sample-and-hold operator; (iv) combined ecology and governance; and (v) a non-periodic or regime-shift null. Freeze training and test windows, transformations, priors, and scoring rules before fitting. Compare held-out log predictive density or another proper score; report phase-order checks and posterior predictive discrepancies, not only in-sample spectral fit.

\textbf{Estimand.} Relative out-of-sample predictive performance of the institutional mechanism versus its alternatives, and the estimated timing/gain parameters conditional on the winning model class.

\textbf{Falsification.} The delay mechanism is weakened when it fails to predict held-out phase ordering or performs worse than an ecological/cohort competitor. It is displaced for that case when the competitor predicts held-out observations better under the pre-specified score. A significant spectrum by itself is not the decision rule.

\subsection{Design 4: randomized human-in-the-loop experiment}

In a laboratory or advisory-interface setting, randomize review cadence and feedback sign or timing in a common renewable-resource simulator. Participants receive a controlled stream of noisy observations and implement decisions through the simulated institution. Randomize treatment assignment, mask the hypothesis where possible, and measure state, command, decision latency, revisions, and errors. A live fishery is not randomized by this proposal.

\textbf{Estimand.} Intention-to-treat effects of assigned review cadence and rule architecture on stability, transient stock loss, command volatility, and decision quality.

\textbf{Falsification.} The proposed gain--phase signature is weakened if randomized cadence does not alter outcomes in the predicted direction or if observed behavioural adaptation eliminates the modeled distinction. Any laboratory result remains externally limited until field validation.

\subsection{Design 5: closed-loop management strategy evaluation}

Construct an MSE with an operating-model ensemble that crosses review interval, decision lag, deployment lag, observation error, rule gain and sign, stock productivity, age structure, environmental forcing, assessment structure, compliance, and process-noise regime. Compare, at minimum, a responsive extractive rule, a responsive protective rule, and a fixed-plan controller; keep catch/effort constraints and objectives identical where meaningful. Use common random numbers across policies and report both local spectral margins and long-run performance metrics.

\textbf{Estimands.} Probability of breaching biological reference points, expected yield and variability, recovery time, implementation burden, distributional indicators, and regret relative to an explicitly stated benchmark.

\textbf{Falsification.} The result does not generalize beyond the MSE if rankings reverse across plausible operating models or if the claimed stability signature does not predict the prespecified performance outcomes. Robustness is reported across model classes, not just parameter draws from one operating model.

\subsection{Cross-cutting preregistration and displacement discipline}

Every design should state whether it tests mechanism, prediction, or policy performance. It should lock data vintage, response windows, rule coding, exclusions, the unit of analysis, multiplicity, uncertainty, and an alternative-model set before outcomes are inspected. The cod case contributes two dated governance events (moratorium announcement 2 July 1992; reopening 26 June 2024 with an 18 kt TAC), but these dates alone do not supply a counterfactual or a governance lead. A fast ecological response does not establish adequacy; annual SSB data do not establish a governance lag.

Three broader hypotheses motivate, but are not tested by, the present work: (a) observation aggregation can change the apparent feedback phase; (b) governance phase ordering can be more diagnostic than periodicity alone; and (c) a substitution certificate can identify when continuous-delay and sampled-data models cease to be interchangeable. Testing them requires the event-linked, high-frequency, independently varied data described above.

